% !TEX root = main.tex
% 正文以图像评价为单位，保留原始实验数值。
\FloatBarrier\Needspace{5\baselineskip}\section{项目背景与检测任务}
\FloatBarrier\Needspace{4\baselineskip}\subsection{应用目标与工作范围}
制造业零件质检需要兼顾缺陷检出、检查效率和设备成本。本项目探索用深度学习识别与边缘计算辅助齿轮外观检查，并连接图像采集、结果显示、自动分拣和数据管理。降低成本、减少漏检和适应复杂环境是设计目标，尚无人工检查基线、完整成本清单或复杂工厂环境对照，不能视为已经验证的效果。

本人负责全部软件与模型工作，成员负责硬件及后续现场部署。报告以算法探索为重点，结合软件如何连接相机、Web页面和STM32控制说明整机方案。目前保留上下两台工业相机，设备已升级为Jetson NX；第一代输送、采集、显示与分拣链路已打通，新模型仅做过简单现场测试，未记录性能数据。

\FloatBarrier\Needspace{4\baselineskip}\subsection{缺陷定义与自建数据}
项目共有50个齿轮：正常10个、划痕15个、侧部缺损15个、齿角缺损10个，均自行采购或定制。采集使用流水线原型的工业摄像头，由人工摆放、调整齿轮姿态获得上下、旋转和侧面等角度，以模拟部分观察条件，尚未覆盖连续输送中的运动和姿态变化。明确可见的划痕和局部齿部缺损均按不合格区域标注，未建立按深度、尺寸或机械性能分级的合格标准。因此检测任务是识别已定义的外观缺陷，不能直接评价齿轮服役安全。

\begin{figure}[H]\centering
\includegraphics[width=\linewidth,height=0.37\textheight,keepaspectratio]{fig01_defect_examples.png}
\caption{自建数据中的正常外观、划痕与齿部缺损示例。红框为人工标注，局部放大用于观察，均不是模型预测。（来源：\protect\hyperlink{source-5}{划痕数据清单}；\protect\hyperlink{source-6}{齿部缺损数据清单}）}
\label{fig:report-1}
\end{figure}

划痕容易与加工纹、高光混淆；齿部缺损的可见性受位置和观察角度影响。标注以当前图像中的可见缺陷为依据。人工确认看不到目标缺陷的图像作为该任务的负样本；可见但难以辨认或准确标注边界的缺陷标为difficult。difficult表示识别困难，不表示缺陷是否存在尚未确定。缺陷工件的某个视角也可能是图像负样本。


\par\medskip\noindent\begin{minipage}{\linewidth}\small\setstretch{1.06}
\captionof{table}{自建图像数据划分}
\begin{tabular}{@{}>{\raggedright\arraybackslash}p{0.230\dimexpr\linewidth-6\tabcolsep\relax}>{\raggedright\arraybackslash}p{0.260\dimexpr\linewidth-6\tabcolsep\relax}>{\raggedright\arraybackslash}p{0.260\dimexpr\linewidth-6\tabcolsep\relax}>{\raggedright\arraybackslash}p{0.250\dimexpr\linewidth-6\tabcolsep\relax}@{}}
\toprule
检测任务 & 训练图像 & 验证图像 & 测试图像 \\
\midrule
划痕 & 364 & 150 & 150 \\
齿部缺损 & 397 & 99 & 150 \\
\bottomrule\end{tabular}\end{minipage}\par\medskip

两类检测使用同一批150张测试图：31张含划痕，44张含齿部缺损，其中4张同时含两类缺陷，即71张至少含一类缺陷、79张无可见目标缺陷。（来源：\protect\hyperlink{source-3}{划痕测试记录}；\protect\hyperlink{source-4}{齿部缺损及联合测试记录}；\protect\hyperlink{source-5}{划痕数据清单}；\protect\hyperlink{source-6}{齿部缺损数据清单}）图像数量不能当作不同工件数量；当前只有图片顺序编号，尚不能确认同一齿轮的多角度图像是否分布在训练、验证和测试中。该限制影响对新工件泛化能力的判断，详细数据检查见\protect\hyperlink{appendix-d}{附录D}。

\FloatBarrier\Needspace{5\baselineskip}\section{技术路线与文献依据}
\FloatBarrier\Needspace{4\baselineskip}\subsection{从图像到检测结果}
系统先定位齿轮并裁剪感兴趣区域（region of interest，ROI），减少背景干扰，再分别检测划痕与齿部缺损。每类缺陷使用两个图像分类模型和一个目标检测模型：前者判断整体外观，后者提供局部缺陷线索，并将最高预测框置信度用于图像级判断。

同类缺陷的三个模型分数经过温度缩放后加权平均，称为分数级加权融合。温度缩放调整分数尺度，使不同模型的分数更适合组合，但当前尚未证明融合分数等于真实缺陷概率。两类融合分数分别与各自阈值比较；只要任一类达到阈值，系统就标记为疑似缺陷。这是基于或规则的决策级融合，详见\protect\hyperlink{appendix-f}{附录F}。

\begin{figure}[H]\centering
\includegraphics[width=\linewidth,height=0.37\textheight,keepaspectratio]{fig02_system_route.png}
\caption{上下双相机、NX、检测软件与分拣机构的关系。虚线为尚待设计的人工复核环节；整机原型连接成功不等于最终模型已完成现场验证。（来源：\protect\hyperlink{project-facts}{项目事实说明}；\protect\href{run:../GearPro_综合技术报告_Markdown含配图/GearPro_综合实验与部署技术报告.md}{《GearPro综合实验与部署技术报告》}）}
\label{fig:report-2}
\end{figure}

软件将同一次检测的分数、阈值、预测类别、预测框和耗时传给页面、统计和串口，避免各处采用不同的判断。历史代码字段PASS/REJECT在本文分别表示“模型未检出缺陷／疑似缺陷”。人工是否复核全部工件、疑似件或抽查模型未检出的工件，仍需确定。仅复核疑似件可以纠正误报，不能自动发现模型漏检。

\FloatBarrier\Needspace{4\baselineskip}\subsection{相关研究如何支持设计}
ResNet和EfficientNet提供两种图像分类骨干\cite{ref10}\cite{ref11}，作为纹理与结构识别的候选；是否适合本项目由齿轮验证数据决定。FPN的多尺度特征思想\cite{ref12}为保留小缺陷细节提供背景，本项目据此比较原始YOLO模型与包含P2检测层的既有配置。P2是较高空间分辨率的检测层，更多细节也可能使正常细纹引起误报。

模型组合已有成熟研究：Kittler等讨论分类器输出的组合规则\cite{ref16}，工业外观检测也采用多检测模型融合\cite{ref17}和按缺陷类别分别建模后的或规则决策融合\cite{ref15}。这些文献说明方法有研究基础，不能据此宣称本项目首创，也不能替代本项目对融合权重的验证。

Guo等的温度缩放研究\cite{ref4}支持对模型分数进行后处理。Saito与Rehmsmeier强调不平衡数据下精确率--召回率分析的价值\cite{ref3}，因此本报告同时展示可见缺陷检出与图像误报。成本敏感学习及Neyman--Pearson分类研究\cite{ref5}\cite{ref6}提示应根据错误代价确定选择目标；本项目只实施了验证集上的经验误报约束，没有实施带总体误差保证的NP算法。

模型评价与选择文献\cite{ref1}\cite{ref2}还提醒：反复用同一验证集选模型、权重和阈值可能产生选择偏差。因此正文区分验证与测试结果，具体参数依据项目记录，论文仅支持方法动机。扩展文献与证据对应见\protect\hyperlink{appendix-a}{附录A}。

\FloatBarrier\Needspace{5\baselineskip}\section{关键发现与最终方案选择}
\FloatBarrier\Needspace{4\baselineskip}\subsection{V3：判断有无缺陷与定位缺陷是不同任务}
\textbf{问题与实验设计。} 系统既要判断图像中有无缺陷，也希望给出局部位置，供查看与复核。V3使用分类模型区分正常、划痕和齿部缺损，另训练目标检测模型定位缺陷，并分别评价这两个任务。齿部缺损的测试记录如下。（来源：\protect\hyperlink{source-11}{V3测试与定位记录}）


\par\medskip\noindent\begin{minipage}{\linewidth}\small\setstretch{1.06}
\captionof{table}{V3图像判断与区域定位的不同评价单位}
\begin{tabular}{@{}>{\raggedright\arraybackslash}p{0.230\dimexpr\linewidth-6\tabcolsep\relax}>{\raggedright\arraybackslash}p{0.260\dimexpr\linewidth-6\tabcolsep\relax}>{\raggedright\arraybackslash}p{0.260\dimexpr\linewidth-6\tabcolsep\relax}>{\raggedright\arraybackslash}p{0.250\dimexpr\linewidth-6\tabcolsep\relax}@{}}
\toprule
评价任务 & 判断对象 & 正确检出／参考数量 & 召回率 \\
\midrule
图像分类 & 图像中是否有齿部缺损 & 15/19张图像 & 78.95\% \\
缺陷定位 & 预测框是否正确对应缺损区域 & 25/52个标注区域 & 48.08\% \\
\bottomrule\end{tabular}\end{minipage}\par\medskip

\textbf{发现与设计意义。} 一张图像可能包含多个缺损区域；分类只需判断这张图有缺陷，定位则需要找到相应区域并满足匹配要求。因此两行分母不同，不能把召回率差值当作分类方法优于检测方法的证据。该实验明确了系统的两项职责：是否标记为疑似缺陷，与能否准确显示每个缺陷位置，应分别验证。

这一评价原则延续到V5。目标检测模型的最高框分数可以参与图像判断，但不代表其预测框已经准确。划痕方案给检测分数0.75的融合系数，表示分数计算中的比例，既不是定位准确率，也不能解释为75\%的检出来自检测模型。软件应按固定分数与阈值作出判定，将预测框作为辅助线索；不能根据有没有显示框改变判断。

V3划痕的整体定位计数与另一套逐框匹配计数存在差异，尚未核清，正文不以其“零检出”作为论据，详细记录见\protect\hyperlink{appendix-d}{附录D}。

\FloatBarrier\Needspace{4\baselineskip}\subsection{V4：排序能力、阈值与集成收益需要分别评价}
\textbf{问题与实验设计。} V4将两类缺陷合并，尝试只判断“有无缺陷”，随后比较不同模型、训练预算和模型组合。下表展示同一既有留出集上的基线、扩展单模型和最终融合结果；具体阈值按验证规则选择，并非在留出集上重新寻找。（来源：\protect\hyperlink{source-12}{V4初始基线记录}；\protect\hyperlink{source-13}{V4最终融合记录}；\protect\hyperlink{source-14}{V4发布配置}；\protect\hyperlink{source-15}{V4候选评价记录}；\protect\hyperlink{source-16}{V4详细技术报告}）


\par\medskip\noindent\begin{minipage}{\linewidth}\small\setstretch{1.06}
\captionof{table}{V4排序指标与所选工作点}
\begin{tabular}{@{}>{\raggedright\arraybackslash}p{0.220\dimexpr\linewidth-8\tabcolsep\relax}>{\raggedright\arraybackslash}p{0.120\dimexpr\linewidth-8\tabcolsep\relax}>{\raggedright\arraybackslash}p{0.240\dimexpr\linewidth-8\tabcolsep\relax}>{\raggedright\arraybackslash}p{0.240\dimexpr\linewidth-8\tabcolsep\relax}>{\raggedright\arraybackslash}p{0.180\dimexpr\linewidth-8\tabcolsep\relax}@{}}
\toprule
方案 & AUPRC & 召回率 & 误报率 & 每图分类前向次数 \\
\midrule
基线模型 & 0.767 & 56.67\%（17/30） & 30.00\%（6/20） & 4 \\
扩展单模型 & 0.899 & 约46.7\% & 约5.0\% & 4 \\
三模型融合 & 0.814 & 63.33\%（19/30） & 35.00\%（7/20） & 12 \\
\bottomrule\end{tabular}\end{minipage}\par\medskip

\ReportNote{表内基线及融合计数核对原始JSON；扩展单模型的留出指标来自既有详细报告中的历史摘要，尚未由逐图预测重算。前向次数按保存的测试时增强配置推导，每个成员执行四次分类。完整实验名、阈值和来源见附录D。}
\textbf{发现一：排序更好，不保证当前阈值下检出更多。} AUPRC综合描述模型在不同分数阈值下的精确率--召回率表现；实际检出和误报则取决于采用哪个阈值。扩展单模型的历史AUPRC较高，但固定判定后的检出更少、误报也更少。因此模型选择与判定阈值选择需要分别说明，不能用单一排序指标替代使用效果，也不能仅凭召回降低断言模型整体变差。当前观察尚不能确定这种变化的唯一原因。

\textbf{发现二：总体检出增加，划痕瓶颈仍在。} 从基线到三模型融合，多检出2张缺陷图像，同时增加1张无可见目标缺陷图像的误报。进一步按缺陷类别核对，划痕始终检出4/11，齿部缺损由13/19升至15/19。因此这次总体提升没有改善划痕检出数，推动后续单独组织划痕数据与模型实验。这里比较的是计数，不能据此认定各模型检出的是完全相同的图像。

\textbf{发现三：模型组合有计算与误报代价。} 最大值融合保留成员中的最高分数，在相同成员分数和阈值下，任一成员高响应即可触发缺陷标记，既可能补检，也可能带入误报。本次基线与融合的成员、阈值均有变化，全部差异不能只归因于最大值规则。三模型各采用四路测试时增强，每图共12次前向；这说明分类计算次数增加，不能写成NX耗时或能耗增加三倍。

基线留出结果曾用于指导扩展方向，因此上述分析用于解释历史选择，不作为始终未查看的盲测证据。它对最终方案的启发是：先查清哪些缺陷没有改善，再决定是否继续增加模型；V5重新比较配置与增强，最终未采用测试时增强。

\FloatBarrier\Needspace{4\baselineskip}\subsection{划痕：组合整体纹理与局部线索}
\textbf{问题与尝试理由。} 针对V4组合后仍未改善的划痕检出，V5重新比较整体纹理与局部线索。分类模型能识别整体纹理，却容易受加工纹影响；目标检测能提供局部线索，但细小划痕的预测框未必准确。实验比较不同模型及其分数组合，检查是否能改善图像级缺陷判断。

\textbf{如何比较。} 划痕实验记录包含6个初步分类模型、2个初步检测模型和2个增加训练轮数上限的分类模型。分类候选包括ResNet18、EfficientNet-B0及YOLO分类器；检测候选比较原始YOLO26n和含P2检测层的模型。模型、融合方案和阈值均依据验证集选择。（来源：\protect\hyperlink{source-1}{划痕发布配置}；\protect\hyperlink{source-7}{划痕候选模型评价表}）

\begin{figure}[H]\centering
\includegraphics[width=\linewidth,height=0.37\textheight,keepaspectratio]{fig05_scratch_candidates.png}
\caption{划痕候选在各自验证阈值下的召回率与误报率，采用当时FPR\ensuremath{\leq}0.20的选择约束。图中是候选配置比较，不能把同时改变的输入尺寸、权重或batch分别解释为单因素效果。（来源：\protect\hyperlink{source-7}{划痕候选模型评价表}）}
\label{fig:report-3}
\end{figure}

\textbf{观察与选择。} 两个检测配置均检出38/43张划痕图像，原始模型误报21/107，P2配置误报19/107，其精确率--召回率曲线下面积（AUPRC）也较高。因此采用P2配置，与两个分类模型组合。最终验证集检出42/43、误报9/107；固定配置后的测试集检出25/31、误报20/119。（来源：\protect\hyperlink{source-1}{划痕发布配置}；\protect\hyperlink{source-3}{划痕测试记录}）


\par\medskip\noindent\begin{minipage}{\linewidth}\small\setstretch{1.06}
\captionof{table}{划痕单模型与融合验证对照}
\begin{tabular}{@{}>{\raggedright\arraybackslash}p{0.380\dimexpr\linewidth-6\tabcolsep\relax}>{\raggedright\arraybackslash}p{0.090\dimexpr\linewidth-6\tabcolsep\relax}>{\raggedright\arraybackslash}p{0.270\dimexpr\linewidth-6\tabcolsep\relax}>{\raggedright\arraybackslash}p{0.260\dimexpr\linewidth-6\tabcolsep\relax}@{}}
\toprule
划痕验证方案 & 成员数 & 正确检出／43张可见划痕图 & 误报／107张无可见划痕图 \\
\midrule
入选EfficientNet单模型 & 1 & 41 & 13 \\
P2检测模型 & 1 & 38 & 19 \\
最终加权融合 & 3 & 42 & 9 \\
\bottomrule\end{tabular}\end{minipage}\par\medskip

\ReportNote{各方案在同一验证集按既有FPR\ensuremath{\leq}0.20规则选择各自阈值。融合比EfficientNet多检出1张、少误报4张，支持当前选型；它不是固定阈值的组件消融，也不是独立测试收益。来源见\protect\hyperlink{source-7}{候选评价表}与\protect\hyperlink{source-1}{发布配置}。}
该结果支持当前配置选择，但两个检测实验的batch不同，尚不能把全部差异归因于P2。目标检测分数可用于图像判断，也不意味着预测框准确；已有框匹配评价中仍存在大量误报。三个模型的错误是否互补，仍需同图错误交集与控制其他条件的消融验证。

\FloatBarrier\Needspace{4\baselineskip}\subsection{齿部缺损：比较标签粒度与困难目标处理}
\textbf{问题与尝试理由。} 齿部缺损在不同位置有不同外观。细分标签可能帮助区分，也可能分散有限样本；困难目标是否保留，会改变训练监督。实验因此同时比较位置标签、困难标注处理及检测结构。

\textbf{如何比较。} 初步检测实验共12组：三类／两类／单类标注，保留／排除困难目标，原始／含P2检测层的YOLO模型，均采用960输入、35轮训练上限。后续还比较分类模型、更大输入和更高的训练轮数上限，完整记录见\protect\hyperlink{appendix-e}{附录E}。（来源：\protect\hyperlink{source-8}{齿部缺损候选模型评价表}）

\begin{figure}[H]\centering
\includegraphics[width=\linewidth,height=0.37\textheight,keepaspectratio]{fig06_missing_matrix.png}
\caption{12组检测配置在验证集上的召回率与误报率。每组采用自身验证阈值，batch等条件未全部一致，因此该矩阵用于配置选择，不是严格隔离每个因素的消融。（来源：\protect\hyperlink{source-8}{齿部缺损候选模型评价表}）}
\label{fig:report-4}
\end{figure}

\textbf{观察与选择。} 单类、排除困难目标的两个结构均检出33/34；原始模型误报6/65，P2配置误报13/65。因此齿部缺损采用单类标注的原始YOLO26n。这里“排除”只作用于训练：删除21个困难标注框，并移除6张仅含困难目标的图像，避免将已知缺陷图误作负样本；验证和测试仍保留完整标注。

最终组合中的分类模型采用不同输入和困难样本设置，检测模型采用初步实验保存的权重。提高训练轮数上限后重新训练的检测模型虽AUPRC更高，但既定选择规则优先比较各位置组中的最低召回率，初步模型约0.917、重新训练模型约0.909，因而前者入选。这一差别涉及少量样本，说明选择过程可核对，不能证明初步实验模型普遍更好。（来源：\protect\href{run:../GearPro_综合技术报告_Markdown含配图/GearPro_综合实验与部署技术报告.md}{《GearPro综合实验与部署技术报告》}；\protect\hyperlink{source-2}{齿部缺损发布配置}）


\par\medskip\noindent\begin{minipage}{\linewidth}\small\setstretch{1.06}
\captionof{table}{齿部缺损单模型与融合验证对照}
\begin{tabular}{@{}>{\raggedright\arraybackslash}p{0.380\dimexpr\linewidth-6\tabcolsep\relax}>{\raggedright\arraybackslash}p{0.090\dimexpr\linewidth-6\tabcolsep\relax}>{\raggedright\arraybackslash}p{0.270\dimexpr\linewidth-6\tabcolsep\relax}>{\raggedright\arraybackslash}p{0.260\dimexpr\linewidth-6\tabcolsep\relax}@{}}
\toprule
齿部缺损验证方案 & 成员数 & 正确检出／34张可见缺损图 & 误报／65张无可见缺损图 \\
\midrule
入选EfficientNet单模型 & 1 & 34 & 6 \\
原始YOLO检测模型 & 1 & 33 & 6 \\
两个分类模型分数均值 & 2 & 34 & 5 \\
EfficientNet与检测分数加权融合 & 2 & 34 & 4 \\
最终加权融合 & 3 & 34 & 4 \\
\bottomrule\end{tabular}\end{minipage}\par\medskip

\ReportNote{各方案在同一验证集按既有FPR\ensuremath{\leq}0.20规则选择各自阈值。融合主要减少误报；两模型与最终三模型计数相同，后者因AUPRC约0.9660高于0.9634入选，不能证明三个成员都必要。完整配置见\protect\hyperlink{appendix-f}{附录F}，来源为\protect\hyperlink{source-9}{融合候选评价表}。}
\FloatBarrier\Needspace{4\baselineskip}\subsection{分别检测是否值得增加模型数量}
两个检测子系统完成后，又训练统一缺陷分类候选，在同一批150张测试图上比较“有无任意缺陷”。分别检测并合并结果的系统检出63/71、误报18/79；已选统一模型检出54/71、误报33/79。前者多检出9张、少误报15张，支持保留当前组合方案。（来源：\protect\hyperlink{source-4}{齿部缺损及联合测试记录}）

两者训练数据构造、模型数量和预算不同，因此比较说明哪套实际系统在当前测试图上更好，尚不能单独证明“分别建模”本身造成全部收益。模型文件也由统一模型约15.6MB增至两个检测子系统约130MB；文件大小不等于NX的运行内存或耗时，部署代价仍需测量。

\FloatBarrier\Needspace{5\baselineskip}\section{参数来源与检测结果}
\FloatBarrier\Needspace{4\baselineskip}\subsection{指标、权重与阈值各自是什么}
召回率、精确率和误报率是通用评价指标，并非YOLO专属参数。为避免混淆，本报告区分学习得到的网络参数、组合分数的融合系数及最终判定阈值。


\begingroup\small\setstretch{1.06}
\begin{longtable}{@{}>{\raggedright\arraybackslash}p{0.240\dimexpr\linewidth-4\tabcolsep\relax}>{\raggedright\arraybackslash}p{0.340\dimexpr\linewidth-4\tabcolsep\relax}>{\raggedright\arraybackslash}p{0.420\dimexpr\linewidth-4\tabcolsep\relax}@{}}
\caption{评价指标与模型参数的区别}\\
\toprule
名称 & 含义或作用 & 来源 \\
\midrule\endfirsthead
\toprule
名称 & 含义或作用 & 来源 \\
\midrule\endhead
\midrule\multicolumn{3}{r}{\footnotesize 续下页}\endfoot
\bottomrule\endlastfoot
召回率（Recall） & 有可见目标缺陷图像中被检出的比例，TP/(TP+FN) & 对照人工标注统计 \\
精确率（Precision） & 被判为有缺陷的图像中真正有缺陷的比例，TP/(TP+FP) & 对照人工标注统计 \\
误报率（FPR） & 无可见目标缺陷图像中被误判的比例，FP/(FP+TN) & 对照人工标注统计 \\
网络参数 & 决定网络如何提取和判断特征 & 预训练初始化后通过训练更新 \\
温度参数 & 调整模型分数尺度 & 选择验证集负对数似然最小的候选 \\
融合系数 & 决定各模型分数的加权比例 & 比较预设融合候选的验证结果 \\
判定阈值 & 将连续分数转换为疑似缺陷标记 & 固定模型与融合后扫描验证阈值 \\
\end{longtable}\endgroup

TP、FN分别为正确检出与漏检的缺陷图像；FP、TN分别为误报与正确排除的无可见目标缺陷图像。指标是评价结果，不直接决定某个权重应该取多少。

\FloatBarrier\Needspace{4\baselineskip}\subsection{核心参数怎样确定}
项目先训练候选模型，再用验证集选择模型、温度、融合和阈值，最后固定配置进行测试。两个分类模型的调整后分数先取平均，再与检测分数加权组合；最终展开系数如下。（来源：\protect\hyperlink{source-1}{划痕发布配置}；\protect\hyperlink{source-2}{齿部缺损发布配置}）


\par\medskip\noindent\begin{minipage}{\linewidth}\small\setstretch{1.06}
\captionof{table}{最终融合系数及判定阈值}
\begin{tabular}{@{}>{\raggedright\arraybackslash}p{0.260\dimexpr\linewidth-8\tabcolsep\relax}>{\raggedright\arraybackslash}p{0.260\dimexpr\linewidth-8\tabcolsep\relax}>{\raggedright\arraybackslash}p{0.160\dimexpr\linewidth-8\tabcolsep\relax}>{\raggedright\arraybackslash}p{0.160\dimexpr\linewidth-8\tabcolsep\relax}>{\raggedright\arraybackslash}p{0.160\dimexpr\linewidth-8\tabcolsep\relax}@{}}
\toprule
检测任务 & EfficientNet分数系数 & ResNet18分数系数 & 检测分数系数 & 判定阈值（约） \\
\midrule
划痕 & 0.125 & 0.125 & 0.750 & 0.3003 \\
齿部缺损 & 0.250 & 0.250 & 0.500 & 0.3413 \\
\bottomrule\end{tabular}\end{minipage}\par\medskip

这些系数来自有限候选比较，并非人为指定的缺陷重要程度，也不是各模型因果贡献。齿部缺损两个候选均得到34TP、4FP，最终方案因AUPRC约0.966高于另一方案约0.963而入选，不能据此宣称全局最优。运行配置保留完整阈值精度，正文仅为阅读取近似值。

历史阈值流程优先满足验证误报率\ensuremath{\leq}0.20，再提高召回率，并结合位置组表现等规则选择。Recall\ensuremath{\geq}0.95、FPR\ensuremath{\leq}0.20缺少正式业务需求与成本测算依据，只作为\textbf{历史实验目标及选择约束}，不作为生产验收标准。 训练损失中的\ensuremath{\gamma}=1.5、齿部缺损正类权重2及ROI边界扩展4\%等属于预先设定的参数，未通过独立实验证明最优；完整来源与选择规则见\protect\hyperlink{appendix-b}{附录B}、\protect\hyperlink{appendix-f}{附录F}。

\FloatBarrier\Needspace{4\baselineskip}\subsection{同一测试集上的结果}
\begin{figure}[H]\centering
\includegraphics[width=\linewidth,height=0.37\textheight,keepaspectratio]{fig08_system_comparison.png}
\caption{同一批150张图像、相同“任意缺陷”标签下的系统比较。误差线为按图像计数的Wilson 95\%区间；工件内图像相关性未知，区间仅作描述性参考，不能据此判断差异具有统计显著性。（来源：\protect\hyperlink{source-4}{齿部缺损及联合测试记录}）}
\label{fig:report-5}
\end{figure}


\par\medskip\noindent\begin{minipage}{\linewidth}\small\setstretch{1.06}
\captionof{table}{固定配置的图像级测试结果}
\begin{tabular}{@{}>{\raggedright\arraybackslash}p{0.260\dimexpr\linewidth-8\tabcolsep\relax}>{\raggedright\arraybackslash}p{0.260\dimexpr\linewidth-8\tabcolsep\relax}>{\raggedright\arraybackslash}p{0.160\dimexpr\linewidth-8\tabcolsep\relax}>{\raggedright\arraybackslash}p{0.160\dimexpr\linewidth-8\tabcolsep\relax}>{\raggedright\arraybackslash}p{0.160\dimexpr\linewidth-8\tabcolsep\relax}@{}}
\toprule
检测方案 & TP/FP/TN/FN & 召回率 & 精确率 & 误报率 \\
\midrule
划痕检测 & 25/20/99/6 & 80.65\% & 55.56\% & 16.81\% \\
齿部缺损检测 & 36/4/102/8 & 81.82\% & 90.00\% & 3.77\% \\
两类结果合并 & 63/18/61/8 & 88.73\% & 77.78\% & 22.78\% \\
统一缺陷分类模型 & 54/33/46/17 & 76.06\% & 62.07\% & 41.77\% \\
\bottomrule\end{tabular}\end{minipage}\par\medskip

前两行分别评价该类缺陷，其负样本中可能有另一类可见缺陷；后两行评价任意缺陷，负样本须无两类可见缺陷。因此只有后两行可直接比较。各类检测的误报约束也不能直接当作合并系统的误报约束。18/79是图像误报计数，不是18个合格工件被误判；缺陷未进入画面属于观察覆盖问题，未计入当前图像级漏检。

验证和测试应分开解释：划痕验证召回为42/43，测试为25/31；齿部缺损验证为34/34，测试为36/44。两项测试召回均约81\%，低于用于选择配置的验证表现。选择偏差、拍摄条件和困难样本覆盖均可能影响差异，目前不能量化各因素贡献。

\FloatBarrier\Needspace{4\baselineskip}\subsection{错误案例与使用取舍}
\begin{figure}[H]\centering
\includegraphics[width=\linewidth,height=0.37\textheight,keepaspectratio]{fig10_scratch_errors.png}
\caption{两个划痕漏检和一个划痕误报案例。红框来自人工标注，模型分数来自固定配置的已有预测记录；局部放大不产生新的推理结果。“无划痕”不意味着无齿部缺损。（来源：\protect\hyperlink{source-10}{划痕逐图预测与标注}）}
\label{fig:report-6}
\end{figure}

浅划痕及正常加工纹是下一步成像和识别的重点。齿部缺损的8张漏检中，7张带有oblique位置标签，提示应重点检查这类缺陷的可见性；不能把标注位置直接等同于采购时的工件类型。

降低阈值可补检，却可能增加误报：历史划痕低阈值约0.2038对应测试召回约93.5\%、误报率约58.8\%。该记录用于解释取舍，不用于在旧测试集上重新挑选阈值。当前合并系统将81/150张图像标为疑似缺陷；尚未建立图像与工件对应及送复核规则，不能把这一比例当作工件送复核比例。研究数据中缺陷占71/150，不能据此估计真实产线工作量，精确率也会随现场缺陷比例变化。\cite{ref3}

\FloatBarrier\Needspace{5\baselineskip}\section{项目贡献与后续验证}
\FloatBarrier\Needspace{4\baselineskip}\subsection{现阶段的项目贡献}
本项目围绕两类外观缺陷分别探索标签、检测结构和分数组合：划痕采用P2配置并组合整体纹理与局部线索，齿部缺损采用单类标注的原始检测结构。验证对照中，划痕融合增加检出并减少误报，齿部缺损融合主要减少误报；这些是当前选型依据，不是新工件上的效果保证。软件将同次判断连接显示、统计与控制输出，形成原型应用链路。

这些工作采用主流网络和既有融合方法。现有证据支持当前图像上的方案选择，尚不足以证明算法首创、生产环境优势或对新齿轮的稳定泛化。尤其是检测框、图像判断和工件分拣应分别评价：图像检出缺陷不代表定位准确，也不代表对应工件已被正确分拣。

\FloatBarrier\Needspace{4\baselineskip}\subsection{下一阶段优先回答的问题}

\par\medskip\noindent\begin{minipage}{\linewidth}\small\setstretch{1.06}
\captionof{table}{下一阶段验证优先级}
\begin{tabular}{@{}>{\raggedright\arraybackslash}p{0.380\dimexpr\linewidth-2\tabcolsep\relax}>{\raggedright\arraybackslash}p{0.620\dimexpr\linewidth-2\tabcolsep\relax}@{}}
\toprule
优先问题 & 后续验证 \\
\midrule
对新工件是否有效 & 建立齿轮编号，按工件与拍摄批次分组，用未参加训练的新齿轮测试 \\
困难缺陷能否拍清 & 控制光照与观察角，记录原图、缺陷可见性和漏检变化 \\
参数与模型组合是否必要 & 控制其他条件比较输入、权重、成员模型与校准，记录共同错误及资源开销 \\
人工复核如何覆盖漏检 & 明确模型未检出工件的复核或抽查方式，测量确认比例和复核时间 \\
NX能否满足工件处理节拍 & 记录采集、定位、识别、传输与分拣对应关系，以及分项耗时和长时运行情况 \\
\bottomrule\end{tabular}\end{minipage}\par\medskip

当前尚未进行部署优化，也没有每件检测耗时记录，不设定未经测量的速度指标。下一阶段先完善工件来源管理、成像、复核和整机测量，再决定是否需要模型优化或更大参数搜索。详细验证建议与软件状态见\protect\hyperlink{appendix-g}{附录G}。

\ReportAnchor{appendix-a}

